\section{结论}

本文建立并审计了一个火星动力下降 SOCP 研究 artifact。其不可替代的工程资产是原始手写 C 源：作者独立、人工维护的 CRS 稀疏矩阵构造及 CRS--CCS 转换实现（\texttt{MarsLanding.c}）。自动生成 C 数据与 Python 路径仅用于一致性交叉验证，不能替代该资产。

历史固定时间手写基线在多条数值路径中给出 400.7 kg 的一位小数目标一致性，但独立审计发现其终端质量低于干重下界约 0.728 kg。因此该数值被保留为 numerical golden，而不作为物理可行燃料结论。

为处理这一问题，本文新增显式干重和区间内下滑锥约束的参数化固定时间 SOCP，并以冻结的自由终端时间网格研究审计 75 条尝试。N=20 无可行 ECOS 候选；N=30、40、60 的最佳已审计 ECOS 点仍随网格变化。只有 N=60 的 ECOS/Clarabel 复算同时通过预注册审计，约为 398.915 kg。该结果表明区间约束和独立审计能够发现节点模型遗漏，但尚不能证明连续时间最优或网格收敛。

工程上，N150 主机的 7.361 ms 仅覆盖 setup 后的 \texttt{ECOS\_solve()} 平均 CPU 时间。本文不支持 p99 或 WCET，也没有目标硬件、内存或端到端证据，因而不主张飞控部署已经完成。下一阶段应先完成同配置 f32/f64 控制实验、目标平台分阶段计时和内存测量，再进行版本化闭环 Monte Carlo、故障注入与安全降级验证。
